\section{Experimental Methodology}
\label{sec:experiments}

% EX-P1 — Research questions
Our evaluation is organized around nine research questions spanning
effectiveness, matched-budget efficiency, optimizer contribution, runtime
behavior, estimator validity, cross-task transfer, heterogeneous evidence,
failure mechanisms, and systems overhead. The logical chain is
straightforward: RQ1 asks whether the system is effective end-to-end; RQ2
whether it is efficient at a fixed budget; RQ3 isolates the optimizer's
contribution from the retrieval substrate; RQ4 examines runtime adaptation;
RQ5 validates the property estimates; RQ6 tests whether the abstraction
survives a change of terminal task; RQ7 whether it generalizes beyond
homogeneous passage evidence; RQ8 attributes failures to a stage; and RQ9
measures systems overhead. Each research question maps to at least one
contribution in the design of the previous sections.

% EX-P2 — Datasets/tasks
The benchmark suite is selected by workload property rather than leaderboard
coverage. Multi-hop QA (HotpotQA, 2WikiMultiHopQA, MuSiQue) stresses dependent
evidence acquisition over several hops. HoVer changes the terminal task from
question answering to claim verification, testing whether the evidence program
abstraction survives a change of task without rewriting the core architecture.
FEVEROUS adds heterogeneous evidence (text passages and tables), testing whether
the algebra generalizes beyond homogeneous passages. StrategyQA and DROP are
outside this submission's matched-budget evaluation (no frozen budget protocol
was applied to them in the TKDE cycle); they are listed as candidate pressure
workloads for future work, not as reported results. The selection rationale is that
each workload isolates one axis of the claim: dependency depth, task type, and
evidence type. Table~\ref{tab:datasets} lists the workloads and their role.

\begin{table}[t]
\centering
\caption{Benchmark workloads and their role in the evaluation. The
SEALED\_TEST main results use the three multi-hop QA sets (8,661 questions
total); HoVer and FEVEROUS are small transfer samples on a different decoder
(\code{qwen3.8-27b}) reported as corroborating evidence. StrategyQA and DROP
are candidate pressure workloads, not reported here.}
\label{tab:datasets}
\small
\begin{tabular}{llllr}
\toprule
workload & task & evidence type & role in paper & $n$ \\
\midrule
HotpotQA & QA & passage & SEALED\_TEST main & 2,863 \\
2WikiMultiHopQA & QA & passage & SEALED\_TEST main & 4,931 \\
MuSiQue & QA & passage & SEALED\_TEST main & 867 \\
HoVer & claim verification & passage & transfer (RQ6) & 15 \\
FEVEROUS & claim verification & text+table & transfer (RQ7) & 2 \\
StrategyQA & QA (implicit) & passage & future work & --- \\
DROP & QA (numerical) & passage & future work & --- \\
\bottomrule
\end{tabular}
\end{table}

% EX-P3 — Controlled substrate
To isolate execution strategy, we control the retrieval substrate wherever
technically possible: corpus version, index, chunking, candidate pool,
generator, answer prompt, and context cap are held fixed across all methods
compared in a cell. The physical-plan optimizer is the only systematic
difference among the arms of an ablation. Baselines are labeled as either
\emph{exact upstream reproductions} (the upstream method's own code, index, and
prompts) or \emph{controlled adaptations} (upstream method adapted to our shared
substrate); the two categories are never conflated in a result table, and an
exact-vs-adapted audit is part of the reproducibility manifest.

% EX-P4 — Baselines
Baselines cover distinct control families rather than a time-ordered list:
single-shot and iterative retrieval; structured planning and executable-program
RAG (PlanRAG, PyRAG); learned adaptive control (DynaKRAG); and static and
cost-only physical planning as diagnostic controls for the proposed optimizer.
In the optimizer ablation, the static arm replays the same frozen logical plan
under the legacy uniform allocation, the cost-only arm optimizes over cost
alone, and the requirement-aware arm applies Equation~\ref{eq:objective}. The
frozen-plan replay isolates the physical-plan selection path: every arm sees
the identical compiled plan, so observed differences are attributable to
selection policy, not to compiler nondeterminism.

% EX-P5 — Metrics
Metrics are reported at three layers. Task metrics measure answer or
verification quality (EM, F1, accuracy). Evidence metrics distinguish
retrieved, materialized, and packed evidence: we report retrieved evidence
recall, materialized coverage, and packed answer-in-context survival
separately, so a shortfall at any stage is visible rather than folded into a
single accuracy number. Systems metrics measure retrieval calls, passages
accessed, reader tokens, LLM invocations, latency, monetary cost where
available, optimizer planning overhead, and provenance completeness. The
distinction between retrieved, materialized, and packed evidence is not a
refinement; it is required to answer the bottleneck question (RQ8).

% EX-P6 — Matched-budget protocol
Primary comparisons are matched by explicit resource constraints rather than by
nominal hyperparameters. The primary matched dimension is realized retrieval
calls: the retrieval-call budget is threaded into the optimizer's search
(\code{retrieval\_budget}), so every arm is compared at the same budget
ceiling, and realized calls are reported per question. We additionally report
quality at common context-token and model-invocation budgets and trace the
realized quality--cost frontier. Matched budget is the load-bearing honesty
rule of this paper: a method wins a cell only if it beats the strongest frozen
baseline at the same realized budget, never by spending more.

Two denominators, reported side by side. Table~\ref{tab:overall} reports exact
match under two denominators. \emph{Answered} EM is computed over the questions
each arm completed successfully (the standard ``accuracy among answered''
view). \emph{All-question} EM (EM$_{\text{all}}$) scores every non-ok question
as EM$=0$---including budget-exceeded and deterministic plan--question
mismatch---so the denominator is the full SEALED\_TEST cell and no failure mode
can be selected away. We report both because the gap between them is itself a
finding: on HotpotQA the chain arm's budget-exceeded count is $323$ versus
$117$ for static, so its answered EM ($56.1\%$) and all-question EM ($54.2\%$)
diverge more than static's ($56.5\%$ vs.\ $54.1\%$), and the all-question
denominator is the conservative one for the deeper-exploring arm. Paired
statistics (bootstrap $\Delta$EM, McNemar) are computed on the shared questions
that both arms answered, because a paired test requires both arms present; the
point estimates of EM$_{\text{all}}$ are reported for transparency and the
paired tests use the answered indicators on shared answered questions.

% EX-P7 — Statistics
Statistical inference follows the paired structure of the evaluation. Primary
hypotheses, metrics, and budgets are specified in the frozen evaluation protocol
(recorded and auditable before any sealed-run execution; see
Section~\ref{sec:experiments}). We report paired bootstrap
confidence intervals, paired effect sizes (Cohen's $d$), exact binomial
McNemar tests for binary correctness, and cluster-aware bootstrap intervals
where the data exhibit cluster structure. Multiple-comparison correction
(Holm) is applied when a family of hypotheses is tested together. A win is
graded: a statistically supported win ($p<0.05$, CI excluding zero, $d>0.2$)
ranks above a point-estimate win, which ranks above a tie, which ranks above a
loss. We do not report a point-estimate difference as a win, and we do not
report $p>0.05$ as equivalence.

% EX-P8 — Reproducibility and contamination control
We maintain explicit development, validation, and sealed-test registries with
fixed deterministic samples (stratified, seed fixed and recorded). All tuning
decisions are logged before sealed evaluation. Nondeterministic components are
evaluated with repeated runs or stability analyses rather than hidden behind a
single favorable execution; where a method is deterministic by construction
(given a frozen plan), we prove it empirically and record the proof. Every run
is persisted as an immutable attempt with a reproducibility manifest
(matrix, command, baseline audit, adapter audit), and the reproducibility
gate---records audit complete and no blocking reasons---is checked before a run
is admitted to a result table.
